\chapter{Container}
\label{cap:container}
\emergencystretch=2.5em

Nel Capitolo 13 abbiamo visto come la virtualizzazione trasformi un server fisico in tante macchine virtuali; nel
capitolo precedente, come la build automation trasformi il codice in un prodotto installabile. Questo capitolo mette
insieme i due temi: i \textbf{container} impacchettano un'applicazione con tutto ciò che le serve per girare, e la
fanno girare allo stesso modo sul portatile dello sviluppatore, sul server di collaudo e nel cloud. Gli esempi sono
stati eseguiti davvero, con Docker 29.9 e Docker Compose 5.6: gli output riportati sono quelli reali.

% Stili delle figure di questo capitolo
\tikzset{
  cbox/.style={rectangle, rounded corners=2pt, draw=#1!70!black, fill=#1!12, line width=0.7pt,
               font=\scriptsize, align=center, inner sep=3pt, minimum height=0.55cm},
  cbox/.default=swblue,
  cstrato/.style={rectangle, draw=#1!75!black, fill=#1!25, line width=0.6pt, font=\scriptsize\bfseries,
                  align=center, minimum height=0.5cm, inner sep=2pt},
  cfreccia/.style={-{Stealth[length=2.2mm]}, line width=0.8pt, draw=swgray},
  ccontainer/.style={rectangle, draw=swblue!80!black, fill=swblue!55, line width=0.7pt,
                     minimum width=#1, minimum height=0.5cm, text=white, font=\scriptsize\bfseries},
  ccontainer/.default=0.9cm,
}

% ============================================================
\section{Il software è cambiato}

L'industria del software è cambiata, e con lei il modo di consegnarlo:

\begin{center}
\begin{tikzpicture}
  \node[card=swgray, text width=5.2cm, align=left, anchor=east] (prima) at (-0.8,0) {%
    \textbf{Prima}\\[2pt]
    \textbullet\ software monolitico\\
    \textbullet\ cicli di sviluppo lunghi\\
    \textbullet\ un solo ambiente di destinazione\\
    \textbullet\ crescita lenta};
  \node[card=swblue, text width=5.2cm, align=left, anchor=west] (ora) at (0.8,0) {%
    \textbf{Ora}\\[2pt]
    \textbullet\ servizi disaccoppiati\\
    \textbullet\ miglioramenti rapidi e iterativi\\
    \textbullet\ molti ambienti di destinazione\\
    \textbullet\ si deve crescere in fretta};
  \draw[cfreccia, line width=2pt] (prima.east) -- (ora.west);
\end{tikzpicture}
\end{center}

E i deployment diventano sempre più complessi:

\begin{itemize}
  \item ogni servizio o componente indipendente usa molti \textbf{stack} tecnologici: linguaggi, framework, database;
  \item ci sono molti \textbf{ambienti} di destinazione diversi: gli ambienti di sviluppo; quelli di pre-produzione, di
        collaudo (\textit{QA}), di \textit{staging}; la produzione, che può essere in sede (\textit{on premises}), in un
        cloud pubblico o in una soluzione ibrida.
\end{itemize}

% ============================================================
\section{La sfida: la «matrice infernale»}

Si immagini un'applicazione fatta di sei servizi, ognuno con il suo stack. Da una parte c'è la molteplicità degli
\textbf{stack}, dall'altra quella degli \textbf{ambienti hardware} su cui devono girare:

\begin{center}
\begin{tikzpicture}[
    srv/.style={cbox=swblue, text width=2.15cm, minimum height=1.25cm},
    amb/.style={cbox=sworange, minimum width=2.2cm}]
  \node[srv] at (0,1.5)    {\textbf{Preprocessing}\\ {\tiny Python 3.10, ffmpeg, libopencv, nodejs, vips, imagemagick}};
  \node[srv] at (2.45,1.5) {\textbf{Analytics}\\ {\tiny OpenJDK 17, Elasticsearch, Grafana}};
  \node[srv] at (4.9,1.5)  {\textbf{Web App}\\ {\tiny Python 3.8, Django, psycopg, libcurl}};
  \node[srv] at (0,0)      {\textbf{App DB}\\ {\tiny PostgreSQL, Postgis, timescale}};
  \node[srv] at (2.45,0)   {\textbf{Queue}\\ {\tiny Redis, redis-sentinel}};
  \node[srv] at (4.9,0)    {\textbf{API Endpoint}\\ {\tiny Python 2.7, Flask, psycopg, gensim}};
  \draw[dashed, line width=0.8pt, draw=swgray] (6.55,2.4) -- (6.55,-0.9);
  \node[font=\Large\bfseries, text=swred] at (6.55,0.75) {?};
  \foreach \y/\t in {1.95/Dev laptop, 1.15/Production server, 0.35/Public cloud, -0.45/QA server} {
    \node[amb] at (8.0,\y) {\t};
  }
  \node[amb] at (10.5,0.75) {Disaster recovery};
  \node[note, align=center] at (2.45,-1.25) {I servizi interagiscono come devono?};
  \node[note, align=center] at (9.2,-1.25) {Si può migrare facilmente da un ambiente all'altro?};
\end{tikzpicture}
\end{center}

\Needspace{14\baselineskip}
Il risultato è la «matrice infernale» (\textit{the Matrix from Hell}): per ogni coppia (servizio, ambiente) bisogna
chiedersi se funzionerà, e con trenta caselle le sorprese sono garantite.

\begin{center}
\renewcommand{\arraystretch}{1.15}
\begin{tabular}{l|ccccc}
 & \textbf{Dev laptop} & \textbf{Production server} & \textbf{Disaster recovery} & \textbf{Public cloud} &
   \textbf{QA server} \\
\hline
Web App & ? & ? & ? & ? & ? \\
API Endpoint & ? & ? & ? & ? & ? \\
Analytics & ? & ? & ? & ? & ? \\
App DB & ? & ? & ? & ? & ? \\
Queue & ? & ? & ? & ? & ? \\
Preprocessing & ? & ? & ? & ? & ? \\
\end{tabular}
\end{center}

% ============================================================
\section{Un'idea dal trasporto merci}

Lo stesso problema l'ha avuto, prima degli anni Sessanta, il \textbf{trasporto merci}. Da una parte una moltitudine di
merci (automobili, pianoforti, televisori, caffè, banane, \ldots), dall'altra una moltitudine di mezzi di trasporto e di
immagazzinamento (treni, camion, gru, carrelli elevatori, navi). Le merci interagiscono fra loro (le banane accanto al
caffè?), e passare da un mezzo all'altro richiede ogni volta di caricare e scaricare tutto, pezzo per pezzo.

La soluzione è stata il \textbf{container}:

\begin{itemize}
  \item è \textbf{standardizzato}: tutti hanno le stesse dimensioni;
  \item può essere caricato con praticamente \textbf{qualunque merce};
  \item viene preparato da chi si occupa della spedizione, che si assicura che all'interno non avvengano interazioni
        indesiderate;
  \item resta \textbf{sigillato} fino alla consegna;
  \item durante il trasporto tutti i container sono uguali: facili da caricare, scaricare, impilare.
\end{itemize}

\begin{center}
\begin{tikzpicture}
  \foreach \y/\t [count=\i] in {1.2/automobili, 0.6/pianoforti, 0/televisori, -0.6/caffè, -1.2/banane}
    \node[cbox=sworange, minimum width=2.0cm] (m\i) at (0,\y) {\t};
  \node[ccontainer=2.2cm, minimum height=1.6cm] (c) at (3.6,0) {container};
  \foreach \y/\t [count=\i] in {1.2/nave, 0.6/treno, 0/camion, -0.6/gru, -1.2/magazzino}
    \node[cbox=swteal, minimum width=2.0cm] (t\i) at (7.2,\y) {\t};
  \foreach \i in {1,...,5} {
    \draw[cfreccia] (m\i.east) -- (c.west);
    \draw[cfreccia] (c.east) -- (t\i.west);
  }
  \node[note] at (0,-1.85) {qualunque merce};
  \node[note] at (3.6,-1.85) {una scatola standard};
  \node[note] at (7.2,-1.85) {qualunque mezzo};
\end{tikzpicture}
\end{center}

% ============================================================
\section{Container per il codice}

L'idea si applica al software così com'è: ogni servizio, con tutto il suo stack, viene chiuso in un
\textbf{container}, e i container si spostano senza cambiamenti su qualunque ambiente.

\begin{center}
\begin{tikzpicture}
  \foreach \y/\t [count=\i] in {1.25/Analytics, 0.75/Preprocessing, 0.25/App DB, -0.25/Web App, -0.75/Queue,
                               -1.25/API Endpoint}
    \node[cbox, minimum width=2.2cm] (s\i) at (0,\y) {\t};
  \node[ccontainer=2.4cm, minimum height=1.6cm] (c) at (3.8,0) {container};
  \foreach \y/\t [count=\i] in {1.2/Dev laptop, 0.6/Production server, 0/Public cloud, -0.6/QA server,
                               -1.2/Disaster recovery}
    \node[cbox=sworange, minimum width=2.6cm] (e\i) at (7.6,\y) {\t};
  \foreach \i in {1,...,6} \draw[cfreccia] (s\i.east) -- (c.west);
  \foreach \i in {1,...,5} \draw[cfreccia] (c.east) -- (e\i.west);
\end{tikzpicture}
\end{center}

Perché darsi la pena? I vantaggi sono diversi per chi sviluppa e per chi gestisce i sistemi:

\begin{center}
\renewcommand{\arraystretch}{1.25}
\begin{tabular}{P{7.4cm}P{7.4cm}}
\toprule
\textbf{Sviluppatori} & \textbf{Operazioni (\textit{operations})} \\
\midrule
devono occuparsi solo di ciò che sta \textbf{dentro} il container &
devono occuparsi solo di ciò che sta \textbf{fuori} dal container \\
l'ambiente di sviluppo si prepara più facilmente &
ogni container si gestisce allo stesso modo \\
niente conflitti fra librerie e dipendenze &
la gestione del ciclo di vita si semplifica \\
\textbf{si costruisce una volta, si esegue ovunque}* &
\textbf{si configura una volta, si esegue qualunque cosa}** \\
\bottomrule
\end{tabular}
\end{center}

\noindent{\footnotesize * ovunque ci siano la stessa architettura e un kernel Linux moderno.
** qualunque cosa costruita per la stessa architettura e lo stesso kernel.}

\medskip
\Needspace{17\baselineskip}
La matrice infernale si risolve così: in ogni casella c'è lo stesso oggetto, un container, che si comporta allo stesso
modo ovunque.

\begin{center}
\renewcommand{\arraystretch}{1.3}
\begin{tabular}{l|ccccc}
 & \textbf{Dev laptop} & \textbf{Production server} & \textbf{Disaster recovery} & \textbf{Public cloud} &
   \textbf{QA server} \\
\hline
Web App & \tikz\node[ccontainer=0.7cm, minimum height=0.32cm]{}; & \tikz\node[ccontainer=0.7cm, minimum height=0.32cm]{}; &
  \tikz\node[ccontainer=0.7cm, minimum height=0.32cm]{}; & \tikz\node[ccontainer=0.7cm, minimum height=0.32cm]{}; &
  \tikz\node[ccontainer=0.7cm, minimum height=0.32cm]{}; \\
API Endpoint & \tikz\node[ccontainer=0.7cm, minimum height=0.32cm]{}; & \tikz\node[ccontainer=0.7cm, minimum height=0.32cm]{}; &
  \tikz\node[ccontainer=0.7cm, minimum height=0.32cm]{}; & \tikz\node[ccontainer=0.7cm, minimum height=0.32cm]{}; &
  \tikz\node[ccontainer=0.7cm, minimum height=0.32cm]{}; \\
Analytics & \tikz\node[ccontainer=0.7cm, minimum height=0.32cm]{}; & \tikz\node[ccontainer=0.7cm, minimum height=0.32cm]{}; &
  \tikz\node[ccontainer=0.7cm, minimum height=0.32cm]{}; & \tikz\node[ccontainer=0.7cm, minimum height=0.32cm]{}; &
  \tikz\node[ccontainer=0.7cm, minimum height=0.32cm]{}; \\
App DB & \tikz\node[ccontainer=0.7cm, minimum height=0.32cm]{}; & \tikz\node[ccontainer=0.7cm, minimum height=0.32cm]{}; &
  \tikz\node[ccontainer=0.7cm, minimum height=0.32cm]{}; & \tikz\node[ccontainer=0.7cm, minimum height=0.32cm]{}; &
  \tikz\node[ccontainer=0.7cm, minimum height=0.32cm]{}; \\
Queue & \tikz\node[ccontainer=0.7cm, minimum height=0.32cm]{}; & \tikz\node[ccontainer=0.7cm, minimum height=0.32cm]{}; &
  \tikz\node[ccontainer=0.7cm, minimum height=0.32cm]{}; & \tikz\node[ccontainer=0.7cm, minimum height=0.32cm]{}; &
  \tikz\node[ccontainer=0.7cm, minimum height=0.32cm]{}; \\
Preprocessing & \tikz\node[ccontainer=0.7cm, minimum height=0.32cm]{}; & \tikz\node[ccontainer=0.7cm, minimum height=0.32cm]{}; &
  \tikz\node[ccontainer=0.7cm, minimum height=0.32cm]{}; & \tikz\node[ccontainer=0.7cm, minimum height=0.32cm]{}; &
  \tikz\node[ccontainer=0.7cm, minimum height=0.32cm]{}; \\
\end{tabular}
\end{center}

% ============================================================
\section{Macchine virtuali e container}

Le macchine virtuali del Capitolo 13 non risolvono gli stessi problemi? Sì, ma non sono altrettanto
\textbf{leggere}: i container permettono deployment, riavvii e spostamenti molto più rapidi. La differenza si vede
negli strati:

\begin{center}
\begin{tikzpicture}[x=1.25cm, y=0.55cm]
  % macchine virtuali
  \node[font=\scriptsize\bfseries] at (1.5,6.0) {Macchine virtuali};
  \node[cstrato=swgray, minimum width=4.6cm] at (1.5,0) {Infrastruttura};
  \node[cstrato=swblue, minimum width=4.6cm] at (1.5,1) {Hypervisor};
  \foreach \x/\a in {0.25/App A, 1.5/App B, 2.75/App B'} {
    \node[cstrato=swteal, minimum width=1.45cm] at (\x,2) {Guest OS};
    \node[cstrato=sworange, minimum width=1.45cm] at (\x,3) {Bins/Libs};
    \node[cstrato=swgreen, minimum width=1.45cm] at (\x,4) {\a};
  }
  % container
  \node[font=\scriptsize\bfseries] at (6.5,6.0) {Container};
  \node[cstrato=swgray, minimum width=4.6cm] at (6.5,0) {Infrastruttura};
  \node[cstrato=swteal, minimum width=4.6cm] at (6.5,1) {Sistema operativo (kernel)};
  \node[cstrato=swblue, minimum width=4.6cm] at (6.5,2) {Container engine};
  \node[cstrato=sworange, minimum width=1.45cm] at (5.25,3) {Bins/Libs};
  \node[cstrato=sworange, minimum width=3.0cm] at (7.125,3) {Bins/Libs};
  \foreach \x/\a in {5.25/App A, 6.5/App B, 7.75/App B'}
    \node[cstrato=swgreen, minimum width=1.45cm] at (\x,4) {\a};
\end{tikzpicture}
\end{center}

Ogni macchina virtuale contiene un intero \textbf{sistema operativo ospite}, con il suo kernel, avviato
dall'hypervisor. I container invece \textbf{condividono il kernel} del sistema che li ospita: ognuno porta con sé solo
le librerie e i programmi dell'applicazione (e due applicazioni possono anche condividerli, come B e B'). Lo si verifica
in un attimo: un container con la distribuzione Alpine Linux, su un host Ubuntu, vede lo stesso kernel dell'host.

\Needspace{11\baselineskip}
\begin{lstlisting}[style=shell, name=Terminale]
$ grep PRETTY_NAME /etc/os-release
PRETTY_NAME="Ubuntu 24.04.5 LTS"
$ uname -r
6.18.44-fc-v230
$ docker run --rm alpine:3.22 grep PRETTY_NAME /etc/os-release
PRETTY_NAME="Alpine Linux v3.22"
$ docker run --rm alpine:3.22 uname -r
6.18.44-fc-v230
\end{lstlisting}

\Needspace{20\baselineskip}
Il file system del container è quello di Alpine, ma il kernel (\texttt{uname -r}) è lo stesso. È il senso
dell'asterisco nella tabella di prima: un container Linux gira ovunque ci sia un kernel Linux compatibile (su Windows e
su macOS, Docker Desktop ne avvia uno in una piccola macchina virtuale). E un container porta con sé il proprio
\textbf{runtime}, indipendente da quello installato sull'host:

\begin{lstlisting}[style=shell, name=Terminale]
$ java -version
openjdk version "21.0.12.1" 2026-08-18
OpenJDK Runtime Environment (build 21.0.12.1+1-1-24.04.4-Ubuntu)
OpenJDK 64-Bit Server VM (build 21.0.12.1+1-1-24.04.4-Ubuntu, mixed mode, sharing)
$ docker run --rm eclipse-temurin:21-jre java -version
Unable to find image 'eclipse-temurin:21-jre' locally
21-jre: Pulling from library/eclipse-temurin
Digest: sha256:0b99b3376fbb50dbf9c0af8535fdd5b8eb0d3fdfc8b1e4967f383b6919f6901a
Status: Downloaded newer image for eclipse-temurin:21-jre
openjdk version "21.0.13" 2026-10-20 LTS
OpenJDK Runtime Environment Temurin-21.0.13+9 (build 21.0.13+9-LTS)
OpenJDK 64-Bit Server VM Temurin-21.0.13+9 (build 21.0.13+9-LTS, mixed mode, sharing)
\end{lstlisting}

Sull'host c'è il Java 21.0.12 di Ubuntu, nel container il 21.0.13 della distribuzione Temurin: un'applicazione nel
container usa sempre il secondo, qualunque cosa sia installata sulla macchina.

% ============================================================
\section{Docker}

\textbf{Docker} (\url{https://www.docker.com/}) è il motore di container (\textit{container engine}) più diffuso:
progetto nato nel 2013, usato da più di 13 milioni di sviluppatori, con più di 9 milioni di applicazioni
«dockerizzate». È lo \textbf{standard de facto} per la containerizzazione del software, ma esistono alternative: LXD,
BuildKit, Buildah, Podman, \ldots

\subsection{L'architettura di Docker}

\begin{center}
\begin{tikzpicture}
  \node[card=swgray, minimum width=3.3cm, minimum height=3.6cm, anchor=north] (cl) at (0,1.8) {};
  \node[font=\scriptsize\bfseries, anchor=north] at (cl.north) {Client};
  \foreach \y/\c in {1.0/docker build, 0.35/docker pull, -0.3/docker run}
    \node[cbox=sworange, font=\scriptsize\ttfamily, minimum width=2.6cm] at (0,\y) {\c};
  \node[note, align=center] at (0,-1.15) {riga di comando,\\ Docker Desktop};
  \node[card=swblue, minimum width=5.4cm, minimum height=3.6cm, anchor=north] (ho) at (4.8,1.8) {};
  \node[font=\scriptsize\bfseries, anchor=north] at (ho.north) {Docker host};
  \node[cbox, minimum width=4.6cm] (d) at (4.8,1.0) {demone Docker (\texttt{dockerd})};
  \node[cbox=swteal, minimum width=2.1cm, minimum height=1.3cm] (im) at (3.65,-0.35) {immagini};
  \node[cbox=swgreen, minimum width=2.1cm, minimum height=1.3cm] (co) at (5.95,-0.35) {container};
  \draw[cfreccia] (im) -- node[font=\tiny, above]{run} (co);
  \node[card=swteal, minimum width=3.4cm, minimum height=3.6cm, anchor=north] (re) at (9.4,1.8) {};
  \node[font=\scriptsize\bfseries, anchor=north] at (re.north) {Registry};
  \node[note, align=center] at (9.4,0.3) {Docker Hub, JFrog,\\ Quay.io, \ldots};
  \draw[cfreccia] (cl.east |- d) -- (d.west);
  \draw[cfreccia] (re.west |- d) -- node[font=\tiny, above]{pull} (d.east);
\end{tikzpicture}
\end{center}

\Needspace{28\baselineskip}
Il \textbf{client} (la riga di comando \texttt{docker}, o l'interfaccia grafica Docker Desktop) manda i comandi al
\textbf{demone} \texttt{dockerd}, che gira sul \textbf{Docker host} e gestisce immagini e container; le immagini che
non ha, il demone le scarica da un \textbf{registry}, come Docker Hub. L'immagine \texttt{hello-world} lo racconta da
sé:

\begin{lstlisting}[style=shell, name=Terminale]
$ docker run --rm hello-world

Hello from Docker!
This message shows that your installation appears to be working correctly.

To generate this message, Docker took the following steps:
 1. The Docker client contacted the Docker daemon.
 2. The Docker daemon pulled the "hello-world" image from the Docker Hub.
    (amd64)
 3. The Docker daemon created a new container from that image which runs the
    executable that produces the output you are currently reading.
 4. The Docker daemon streamed that output to the Docker client, which sent it
    to your terminal.

To try something more ambitious, you can run an Ubuntu container with:
 $ docker run -it ubuntu bash

Share images, automate workflows, and more with a free Docker ID:
 https://hub.docker.com/

For more examples and ideas, visit:
 https://docs.docker.com/get-started/
\end{lstlisting}

L'opzione \texttt{--rm} cancella il container appena il programma termina.

\subsection{Le immagini}

\dfn{Immagine Docker}{
  Un'\textbf{immagine} è un modello (\textit{template}) \textbf{portabile} e \textbf{di sola lettura}, che contiene
  tutte le istruzioni per creare un container. Un \textbf{container} è un'istanza in esecuzione di un'immagine, come un
  oggetto lo è di una classe.}

\Needspace{13\baselineskip}
Un'immagine si può ottenere in tre modi:

\begin{center}
\begin{tikzpicture}
  \node[cbox=swgray, minimum width=2.4cm] (ar) at (0,0) {archivio \texttt{.tar}};
  \node[cbox=sworange, minimum width=2.4cm] (df) at (3.6,0.6) {Dockerfile};
  \node[cbox=swteal, minimum width=2.4cm] (rg) at (7.2,0) {registry};
  \node[cbox, minimum width=2.6cm, minimum height=0.8cm] (im) at (3.6,-1.4) {\textbf{immagine}};
  \node[ccontainer=2.6cm, minimum height=0.8cm] (co) at (3.6,-3.0) {container};
  \draw[cfreccia] (ar) -- node[font=\tiny\ttfamily, below left]{load} (im);
  \draw[cfreccia] (df) -- node[font=\tiny\ttfamily, right]{build} (im);
  \draw[cfreccia] (rg) -- node[font=\tiny\ttfamily, below right]{pull} (im);
  \draw[cfreccia] (im) -- node[font=\tiny\ttfamily, right]{run} (co);
\end{tikzpicture}
\end{center}

\begin{itemize}
  \item si \textbf{carica} da un archivio \texttt{tar} (\texttt{docker load}), creato per esempio con
        \texttt{docker save};
  \item si \textbf{scarica} da un registry (\texttt{docker pull});
  \item si \textbf{costruisce} estendendo un'immagine esistente con un elenco di istruzioni scritte in un file di testo,
        il \textbf{Dockerfile} (\texttt{docker build}).
\end{itemize}

\Needspace{21\baselineskip}
Le prime due, all'opera:

\begin{lstlisting}[style=shell, name=Terminale]
$ docker pull postgres:17-alpine
17-alpine: Pulling from library/postgres
dddde5442af4: Pulling fs layer
e018eb860c5e: Pulling fs layer
...
d6e09b1bfe5f: Pull complete
Digest: sha256:266d3c257594e195b309ae9d1805e54c788f1d7b375f56bf6c367aac6c917e7b
Status: Downloaded newer image for postgres:17-alpine
docker.io/library/postgres:17-alpine
$ docker save alpine:3.22 -o alpine.tar
$ du -h alpine.tar
3.8M	alpine.tar
$ docker rmi alpine:3.22
Untagged: alpine:3.22
Deleted: sha256:2e0b81d60597f449baa56dc25d355367869d2dca4cd57e866f1e2b52af3eb55f
$ docker load -i alpine.tar
Loaded image: alpine:3.22
\end{lstlisting}

Un'immagine è fatta di \textbf{strati} (\textit{layer}): \texttt{docker pull} ne scarica una decina per PostgreSQL. Il
\textbf{nome} di un'immagine ha la forma \texttt{nome:etichetta} (\texttt{postgres:17-alpine}: PostgreSQL 17, costruito
su Alpine), e l'impronta \texttt{sha256} ne identifica esattamente il contenuto. L'immagine di Alpine, salvata in un
archivio, occupa meno di 4 MB: è uno dei motivi per cui i container sono leggeri.

% ============================================================
\Needspace{56\baselineskip}
\section{Un esempio: un servizio Spring Boot}

L'esempio delle slide è un microservizio, \texttt{servizi-tasse}, che calcola tasse e agevolazioni per l'«Ateneo
Futura». È scritto con \textbf{Spring Boot}, un framework Java per applicazioni web, e costruito con Maven (capitolo
precedente). Ecco il suo POM, dalle slide:

\begin{lstlisting}[style=xml, name=servizi-tasse/pom.xml]
<?xml version="1.0" encoding="UTF-8"?>
<project xmlns="http://maven.apache.org/POM/4.0.0"
     xmlns:xsi="http://www.w3.org/2001/XMLSchema-instance"
     xsi:schemaLocation="http://maven.apache.org/POM/4.0.0
https://maven.apache.org/xsd/maven-4.0.0.xsd">
  <modelVersion>4.0.0</modelVersion>
  <parent>
    <groupId>org.springframework.boot</groupId>
    <artifactId>spring-boot-starter-parent</artifactId>
    <version>3.2.3</version>
    <relativePath/>
  </parent>
  <groupId>it.unifutura</groupId>
  <artifactId>servizi-tasse</artifactId>
  <version>1.0.0</version>
  <name>servizi-tasse</name>
  <description>Microservizi calcolo tasse e agevolazioni per Ateneo
Futura</description>
  <properties>
    <java.version>21</java.version>
    <springdoc.version>2.3.0</springdoc.version>
  </properties>
  <dependencies>
    <dependency>
       <groupId>org.springframework.boot</groupId>
       <artifactId>spring-boot-starter-web</artifactId>
    </dependency>
    <dependency>
      <groupId>org.springframework.boot</groupId>
      <artifactId>spring-boot-starter-validation</artifactId>
    </dependency>
    <dependency>
      <groupId>org.springdoc</groupId>
      <artifactId>springdoc-openapi-starter-webmvc-ui</artifactId>
      <version>${springdoc.version}</version>
    </dependency>
    <dependency>
      <groupId>org.springframework.boot</groupId>
      <artifactId>spring-boot-starter-test</artifactId>
      <scope>test</scope>
    </dependency>
  </dependencies>
  <build>
    <plugins>
      <plugin>
        <groupId>org.springframework.boot</groupId>
        <artifactId>spring-boot-maven-plugin</artifactId>
      </plugin>
    </plugins>
  </build>
</project>
\end{lstlisting}

Il POM \textbf{eredita} da un POM \texttt{parent}, \texttt{spring-boot-starter-parent}, che fissa le versioni di tutte
le librerie di Spring: per questo le dipendenze non indicano la versione, tranne \texttt{springdoc}, che di Spring non
fa parte (genera la documentazione OpenAPI del servizio, Capitolo 14). Gli \textit{starter} sono dipendenze che ne
raccolgono molte altre: \texttt{web} porta con sé anche il server web Tomcat. Il plugin di Spring Boot, infine,
trasforma il \texttt{jar} in un \textbf{eseguibile autosufficiente}, con tutte le dipendenze dentro.

Il codice del servizio non è nelle slide: qui ne usiamo uno minimo, con regole inventate per l'esempio.

\Needspace{14\baselineskip}
\begin{lstlisting}[style=java, name=ServiziTasseApplication.java]
package it.unifutura.servizitasse;

import org.springframework.boot.SpringApplication;
import org.springframework.boot.autoconfigure.SpringBootApplication;

@SpringBootApplication
public class ServiziTasseApplication {
    public static void main(String[] args) {
        SpringApplication.run(ServiziTasseApplication.class, args);
    }
}
\end{lstlisting}

\Needspace{21\baselineskip}
\begin{lstlisting}[style=java, name=TasseController.java]
package it.unifutura.servizitasse;

import org.springframework.web.bind.annotation.GetMapping;
import org.springframework.web.bind.annotation.RequestParam;
import org.springframework.web.bind.annotation.RestController;

@RestController
public class TasseController {

    record Tassa(int isee, String fascia, int importo) {}

    // regole inventate per l'esempio: esonero fino a 22.000 euro di ISEE
    @GetMapping("/tasse")
    public Tassa calcola(@RequestParam int isee) {
        if (isee <= 22000) return new Tassa(isee, "esonero", 0);
        if (isee <= 40000) return new Tassa(isee, "intermedia", 900);
        return new Tassa(isee, "massima", 2000);
    }
}
\end{lstlisting}

\Needspace{18\baselineskip}
\texttt{@RestController} e \texttt{@GetMapping} dicono a Spring che il metodo risponde alle richieste \texttt{GET} per
\texttt{/tasse}, e il record restituito viene convertito in JSON (è un servizio REST, Capitolo 14). Con Maven si ottiene
il jar:

\begin{lstlisting}[style=shell, name=Terminale]
$ mvn -ntp clean package
...
[INFO] Building jar: [cartella del progetto]/target/servizi-tasse-1.0.0.jar
[INFO]
[INFO] --- spring-boot:3.2.3:repackage (repackage) @ servizi-tasse ---
[INFO] Replacing main artifact [cartella del progetto]/target/servizi-tasse-1.0.0.jar with repackaged archive, adding nested dependencies in BOOT-INF/.
[INFO] ------------------------------------------------------------------------
[INFO] BUILD SUCCESS
...
$ du -h target/servizi-tasse-1.0.0.jar
26M	target/servizi-tasse-1.0.0.jar
\end{lstlisting}

(L'opzione \texttt{-ntp} nasconde l'avanzamento dei download.) 26 MB, perché dentro ci sono Spring, Tomcat e tutte le
altre dipendenze.

\Needspace{15\baselineskip}
\subsection{Il Dockerfile}

L'immagine si descrive con un \textbf{Dockerfile}:

\begin{lstlisting}[style=docker, name=servizi-tasse/Dockerfile]
# si parte da un'immagine che contiene solo il runtime di Java 21
FROM eclipse-temurin:21-jre
WORKDIR /app
# il jar prodotto da "mvn package", con tutte le dipendenze dentro
COPY target/servizi-tasse-1.0.0.jar app.jar
EXPOSE 8080
ENTRYPOINT ["java", "-jar", "app.jar"]
\end{lstlisting}

\begin{itemize}
  \item \texttt{FROM} sceglie l'immagine di partenza, che si estende: qui una con il solo runtime di Java (JRE, senza il
        compilatore);
  \item \texttt{WORKDIR} imposta la cartella di lavoro dentro l'immagine;
  \item \texttt{COPY} copia un file dalla cartella del progetto (il \textbf{contesto} della build) nell'immagine;
  \item \texttt{EXPOSE} documenta la porta su cui il servizio ascolta;
  \item \texttt{ENTRYPOINT} è il comando eseguito all'avvio del container.
\end{itemize}

\Needspace{33\baselineskip}
Ogni istruzione produce uno strato dell'immagine. Si costruisce con \texttt{docker build} e si avvia con
\texttt{docker run}:

\begin{lstlisting}[style=shell, name=Terminale]
$ docker build -t servizi-tasse:1.0 .
...
#5 [1/3] FROM docker.io/library/eclipse-temurin:21-jre@sha256:0b99b3376fbb50dbf9c0af8535fdd5b8eb0d3fdfc8b1e4967f383b6919f6901a
#5 resolve docker.io/library/eclipse-temurin:21-jre@sha256:0b99b3376fbb50dbf9c0af8535fdd5b8eb0d3fdfc8b1e4967f383b6919f6901a done
#5 DONE 0.0s

#6 [2/3] WORKDIR /app
#6 CACHED

#7 [3/3] COPY target/servizi-tasse-1.0.0.jar app.jar
#7 CACHED
...
#8 naming to docker.io/library/servizi-tasse:1.0 done
#8 unpacking to docker.io/library/servizi-tasse:1.0 done
#8 DONE 0.1s
$ docker images servizi-tasse
IMAGE               ID             DISK USAGE   CONTENT SIZE   EXTRA
servizi-tasse:1.0   5f446832236f        509MB          140MB
$ docker run -d --name tasse -p 8090:8080 servizi-tasse:1.0
10adc5ebeea14f9f434c190da07d40d589f2984c707556d3f6767f17c0de8e86
$ docker ps
CONTAINER ID   IMAGE                         COMMAND               CREATED         STATUS         PORTS                    NAMES
10adc5ebeea1   servizi-tasse:1.0             "java -jar app.jar"   3 seconds ago   Up 2 seconds   0.0.0.0:8090->8080/tcp   tasse
3c16f233d104   servizi-tasse-servizi-tasse   "java -jar app.jar"   6 minutes ago   Up 5 minutes   0.0.0.0:8080->8080/tcp   unifutura-servizi-tasse
$ curl -s 'localhost:8090/tasse?isee=30000'; echo
{"isee":30000,"fascia":"intermedia","importo":900}
$ docker stop tasse
tasse
$ docker rm tasse
tasse
\end{lstlisting}

I passi della build segnati \texttt{CACHED} non sono stati rifatti: Docker riusa gli strati già costruiti, come Make e
Gradle riusano i file già prodotti. \texttt{-t} dà il nome all'immagine; \texttt{docker run -d} avvia il container
in background (e ne stampa l'identificativo), \texttt{--name} gli dà un nome, \texttt{-p 8090:8080} collega la porta
8090 dell'host alla porta 8080 del container. \texttt{docker ps} elenca i container in esecuzione: il secondo è quello
avviato con Docker Compose, che vediamo fra poco. \texttt{docker stop} ferma il container, \texttt{docker rm} lo
cancella.

\Needspace{17\baselineskip}
\subsection{Docker Compose}

Con più opzioni, o più container, la riga di comando diventa scomoda: \textbf{Docker Compose} descrive l'intera
applicazione in un file YAML, \texttt{docker-compose.yml}. Quello delle slide:

\begin{lstlisting}[style=yaml, name=servizi-tasse/docker-compose.yml]
version: '3.8'
services:
 servizi-tasse:
  build: .
  container_name: unifutura-servizi-tasse
  ports:
   - "8080:8080"
  restart: always
\end{lstlisting}

\begin{itemize}
  \item c'è un solo \textbf{servizio}, \texttt{servizi-tasse};
  \item la sua immagine si costruisce (\texttt{build}) dalla cartella corrente, cioè con il Dockerfile che vi si trova;
  \item il container si chiama \texttt{unifutura-servizi-tasse};
  \item è raggiungibile su \texttt{localhost}, alla porta 8080;
  \item viene riavviato sempre (\texttt{restart: always}), anche a ogni riavvio di Docker.
\end{itemize}

\Needspace{27\baselineskip}
\texttt{docker compose up --build} controlla (o scarica) le risorse necessarie, costruisce le immagini, crea i
container e li avvia; \texttt{-d} lo fa in background:

\begin{lstlisting}[style=shell, name=Terminale]
$ docker compose up --build -d
time="2026-11-03T19:11:04Z" level=warning msg="[cartella del progetto]/docker-compose.yml: the attribute `version` is obsolete, it will be ignored, please remove it to avoid potential confusion"
 Image servizi-tasse-servizi-tasse Building
...
#8 [3/3] COPY target/servizi-tasse-1.0.0.jar app.jar
#8 DONE 0.1s
...
 Image servizi-tasse-servizi-tasse Built
 Network servizi-tasse_default Creating
 Network servizi-tasse_default Creating
 Network servizi-tasse_default Created
 Network servizi-tasse_default Created
 Container unifutura-servizi-tasse Creating
 Container unifutura-servizi-tasse Created
 Container unifutura-servizi-tasse Starting
 Container unifutura-servizi-tasse Started
$ curl -s 'localhost:8080/tasse?isee=15000'; echo
{"isee":15000,"fascia":"esonero","importo":0}
$ curl -s 'localhost:8080/tasse?isee=35000'; echo
{"isee":35000,"fascia":"intermedia","importo":900}
$ curl -s 'localhost:8080/tasse?isee=60000'; echo
{"isee":60000,"fascia":"massima","importo":2000}
\end{lstlisting}

(Alcune righe sono ripetute: è l'output di Compose quando non scrive su un terminale interattivo, dove invece mostra un
riepilogo che si aggiorna sul posto.) Compose ha costruito l'immagine, chiamandola con il nome del progetto (quello della cartella)
seguito da quello del servizio, ha creato una \textbf{rete} per i container del progetto e ha avviato il container.

\warning{L'attributo \texttt{version} è obsoleto}{
  L'avviso iniziale riguarda la prima riga del file: nelle versioni attuali di Compose l'attributo \texttt{version}
  non serve più, e viene ignorato. Basta toglierlo, come faremo nell'esempio successivo.}

\warning{Una svista nelle slide: \texttt{mvn build}}{
  Nelle slide il servizio è «costruito per impostazione predefinita chiamando \texttt{mvn build}». Ma \texttt{build}
  non è una fase di Maven, e \texttt{mvn build} termina con l'errore \texttt{Unknown lifecycle phase "build"} (capitolo
  precedente); né Compose chiama Maven da sé: \texttt{build: .} costruisce l'immagine seguendo il Dockerfile. Maven va
  eseguito prima, come qui (\texttt{mvn package}), oppure dentro il Dockerfile stesso, con una build in più stadi
  (\textit{multi-stage}): un primo stadio basato su un'immagine con Maven compila il jar, e il secondo, con il solo
  runtime, lo copia dal primo.}

\Needspace{19\baselineskip}
I log del container mostrano Spring Boot che si avvia:

\begin{lstlisting}[style=shell, name=Terminale]
$ docker compose logs --no-log-prefix servizi-tasse
...
  .   ____          _            __ _ _
 /\\ / ___'_ __ _ _(_)_ __  __ _ \ \ \ \
( ( )\___ | '_ | '_| | '_ \/ _` | \ \ \ \
 \\/  ___)| |_)| | | | | || (_| |  ) ) ) )
  '  |____| .__|_| |_|_| |_\__, | / / / /
 =========|_|==============|___/=/_/_/_/
 :: Spring Boot ::                (v3.2.3)

2026-11-03T19:11:07.513Z  INFO 1 --- [           main] i.u.s.ServiziTasseApplication            : Starting ServiziTasseApplication v1.0.0 using Java 21.0.13 with PID 1 (/app/app.jar started by root in /app)
...
2026-11-03T19:11:09.381Z  INFO 1 --- [           main] o.s.b.w.embedded.tomcat.TomcatWebServer  : Tomcat started on port 8080 (http) with context path ''
2026-11-03T19:11:09.397Z  INFO 1 --- [           main] i.u.s.ServiziTasseApplication            : Started ServiziTasseApplication in 2.32 seconds (process running for 2.728)
...
\end{lstlisting}

\Needspace{15\baselineskip}
Due dettagli: il Java è il 21.0.13 dell'immagine, e il programma ha \textbf{PID 1}. Dentro il container,
l'applicazione è il primo e unico processo: vede solo sé stessa. Se termina, termina anche il container; e qui
entra in gioco \texttt{restart: always}. Simuliamo un crash, mandando al processo 1 il segnale di terminazione:

\begin{lstlisting}[style=shell, name=Terminale]
$ docker exec unifutura-servizi-tasse kill 1
$ docker inspect -f '{{.RestartCount}}' unifutura-servizi-tasse
1
$ docker ps --format 'table {{.Names}}\t{{.Status}}'
NAMES                     STATUS
unifutura-servizi-tasse   Up 2 seconds
$ curl -s 'localhost:8080/tasse?isee=15000'; echo
{"isee":15000,"fascia":"esonero","importo":0}
\end{lstlisting}

\texttt{docker exec} esegue un comando dentro un container in funzione. Il container si è fermato ed è stato
riavviato da Docker (\texttt{RestartCount} vale 1, ed è «su» da due secondi), e il servizio risponde di nuovo; nei log
compare un secondo avvio di Spring Boot.

% ============================================================
\section{Due container che collaborano}

Il secondo esempio delle slide ha due servizi: uno gestisce gli utenti (\texttt{st-user-service}), l'altro
l'autenticazione (\texttt{st-auth-service}), e il secondo interroga il primo.

\Needspace{23\baselineskip}
\begin{lstlisting}[style=yaml, name=st-microservices/docker-compose.yml]
version: '3.8'

services:
 st-user-service:
  build:
    context: ./st-user-service
    dockerfile: Dockerfile
  container_name: st-user-service
  ports:
    - "0.0.0.0:8082:8082"

 st-auth-service:
  build:
    context: ./st-auth-service
    dockerfile: Dockerfile
  container_name: st-auth-service
  ports:
    - "0.0.0.0:8081:8081"
  environment:
    - ST_USER_SERVICE_URL=http://st-user-service:8082
  depends_on:
    - st-user-service
\end{lstlisting}

\begin{itemize}
  \item ogni servizio ha la sua cartella (\texttt{context}), con il suo Dockerfile e il suo \texttt{pom.xml};
  \item \texttt{0.0.0.0:8082:8082} pubblica la porta su tutte le interfacce di rete dell'host;
  \item \texttt{environment} passa al servizio di autenticazione una \textbf{variabile d'ambiente} con l'indirizzo
        dell'altro servizio: il codice non contiene indirizzi scritti a mano;
  \item \texttt{depends\_on} fa avviare \texttt{st-user-service} prima di \texttt{st-auth-service}. Attenzione: decide
        solo l'\textbf{ordine di avvio}, non aspetta che l'altro servizio sia pronto a rispondere.
\end{itemize}

\Needspace{35\baselineskip}
I due POM, dalle slide, sono quasi identici: cambia l'\texttt{artifactId}. Quello del primo servizio (qui completato con
la dipendenza \texttt{web} e il plugin di Spring Boot, che nelle slide sono omessi):

\begin{lstlisting}[style=xml, name=st-user-service/pom.xml]
<?xml version="1.0" encoding="UTF-8"?>
<project xmlns="http://maven.apache.org/POM/4.0.0"
      xmlns:xsi="http://www.w3.org/2001/XMLSchema-instance"
     xsi:schemaLocation="http://maven.apache.org/POM/4.0.0
https://maven.apache.org/xsd/maven-4.0.0.xsd">
    <modelVersion>4.0.0</modelVersion>
    <parent>
      <groupId>org.springframework.boot</groupId>
      <artifactId>spring-boot-starter-parent</artifactId>
      <version>3.2.3</version>
    </parent>
    <groupId>com.edu.microservices</groupId>
    <artifactId>st-user-service</artifactId>
    <version>1.0.0</version>
    <properties>
      <java.version>17</java.version>
    </properties>
    <dependencies>
      <dependency>
        <groupId>org.springframework.boot</groupId>
        <artifactId>spring-boot-starter-web</artifactId>
      </dependency>
    </dependencies>
    <build>
      <plugins>
        <plugin>
          <groupId>org.springframework.boot</groupId>
          <artifactId>spring-boot-maven-plugin</artifactId>
        </plugin>
      </plugins>
    </build>
</project>
\end{lstlisting}

Anche qui il codice è nostro. Il servizio degli utenti tiene i dati in memoria e li restituisce per nome utente;
ascolta sulla porta 8082 (lo dice il file \texttt{application.properties}, con \texttt{server.port=8082}).

\Needspace{31\baselineskip}
\begin{lstlisting}[style=java, name=UserServiceApplication.java]
package com.edu.microservices;

import java.util.Map;
import org.springframework.boot.SpringApplication;
import org.springframework.boot.autoconfigure.SpringBootApplication;
import org.springframework.http.ResponseEntity;
import org.springframework.web.bind.annotation.GetMapping;
import org.springframework.web.bind.annotation.PathVariable;
import org.springframework.web.bind.annotation.RestController;

@SpringBootApplication
@RestController
public class UserServiceApplication {

    record Utente(String username, String nome, String ruolo) {}

    // i dati degli utenti: in un servizio vero starebbero in un database
    private final Map<String, Utente> utenti = Map.of(
        "janet", new Utente("janet", "Janet Weiss", "studente"),
        "frank", new Utente("frank", "Frank Furter", "docente"));

    @GetMapping("/users/{username}")
    public ResponseEntity<Utente> utente(@PathVariable String username) {
        Utente u = utenti.get(username);
        return u == null ? ResponseEntity.notFound().build() : ResponseEntity.ok(u);
    }

    public static void main(String[] args) {
        SpringApplication.run(UserServiceApplication.class, args);
    }
}
\end{lstlisting}

\Needspace{41\baselineskip}
Il servizio di autenticazione (porta 8081) legge l'indirizzo dell'altro dalla variabile d'ambiente
\texttt{ST\_USER\_SERVICE\_URL} e lo interroga con un client HTTP (\texttt{RestClient}):

\begin{lstlisting}[style=java, name=AuthServiceApplication.java]
package com.edu.microservices;

import java.util.Map;
import org.springframework.beans.factory.annotation.Value;
import org.springframework.boot.SpringApplication;
import org.springframework.boot.autoconfigure.SpringBootApplication;
import org.springframework.web.bind.annotation.GetMapping;
import org.springframework.web.bind.annotation.PathVariable;
import org.springframework.web.bind.annotation.RestController;
import org.springframework.web.client.HttpClientErrorException;
import org.springframework.web.client.RestClient;

@SpringBootApplication
@RestController
public class AuthServiceApplication {

    private final RestClient userService;

    // l'indirizzo dell'altro servizio arriva da una variabile d'ambiente
    public AuthServiceApplication(@Value("${ST_USER_SERVICE_URL}") String url) {
        this.userService = RestClient.create(url);
    }

    @GetMapping("/auth/{username}")
    public Map<String, Object> verifica(@PathVariable String username) {
        try {
            Map<?, ?> utente = userService.get().uri("/users/{u}", username)
                                          .retrieve().body(Map.class);
            return Map.of("username", username, "autorizzato", true,
                          "ruolo", utente.get("ruolo"));
        } catch (HttpClientErrorException.NotFound e) {
            return Map.of("username", username, "autorizzato", false);
        }
    }

    public static void main(String[] args) {
        SpringApplication.run(AuthServiceApplication.class, args);
    }
}
\end{lstlisting}

\Needspace{10\baselineskip}
I Dockerfile sono come quello di prima, con il runtime di Java 17:

\begin{lstlisting}[style=docker, name=st-auth-service/Dockerfile]
FROM eclipse-temurin:17-jre
WORKDIR /app
COPY target/st-auth-service-1.0.0.jar app.jar
EXPOSE 8081
ENTRYPOINT ["java", "-jar", "app.jar"]
\end{lstlisting}

\Needspace{30\baselineskip}
Compilati i due servizi con \texttt{mvn package}, un solo comando costruisce e avvia tutto:

\begin{lstlisting}[style=shell, name=Terminale]
$ docker compose up --build -d
time="2026-11-03T19:12:16Z" level=warning msg="[cartella del progetto]/docker-compose.yml: the attribute `version` is obsolete, it will be ignored, please remove it to avoid potential confusion"
 Image st-microservices-st-user-service Building
 Image st-microservices-st-auth-service Building
...
 Image st-microservices-st-user-service Built
 Image st-microservices-st-auth-service Built
...
 Network st-microservices_default Created
...
 Container st-user-service Created
...
 Container st-auth-service Created
...
 Container st-user-service Started
...
 Container st-auth-service Started
$ curl -s localhost:8081/auth/janet; echo
{"username":"janet","ruolo":"studente","autorizzato":true}
$ curl -s localhost:8081/auth/frank; echo
{"username":"frank","ruolo":"docente","autorizzato":true}
$ curl -s localhost:8081/auth/riff; echo
{"username":"riff","autorizzato":false}
$ curl -s localhost:8082/users/janet; echo
{"username":"janet","nome":"Janet Weiss","ruolo":"studente"}
$ curl -s -o /dev/null -w '%{http_code}\n' localhost:8082/users/riff
404
\end{lstlisting}

Le immagini sono costruite in parallelo, e i container avviati nell'ordine voluto da \texttt{depends\_on}. Una
richiesta a \texttt{st-auth-service} produce una richiesta a \texttt{st-user-service}: per \texttt{riff}, che non
esiste, il servizio degli utenti risponde \texttt{404} e quello di autenticazione nega l'accesso.

\Needspace{17\baselineskip}
\subsection{Come si trovano i container}

L'indirizzo \texttt{http://st-user-service:8082} usa il \textbf{nome} del container come nome di host. Funziona perché
Compose ha creato una \textbf{rete} per il progetto, e dentro quella rete Docker fa da DNS: ogni container è
raggiungibile con il suo nome. Fuori dalla rete, invece, quel nome non significa nulla:

\begin{lstlisting}[style=shell, name=Terminale]
$ docker network ls --filter name=st-microservices
NETWORK ID     NAME                       DRIVER    SCOPE
2a0f62ae1107   st-microservices_default   bridge    local
$ docker exec st-auth-service getent hosts st-user-service
172.19.0.2      st-user-service
$ getent hosts st-user-service || echo 'nome sconosciuto fuori dalla rete di Compose'
nome sconosciuto fuori dalla rete di Compose
\end{lstlisting}

Dentro il container di autenticazione il nome si risolve nell'indirizzo privato dell'altro container (172.19.0.2);
sull'host no. La rete è di tipo \texttt{bridge}: una rete virtuale, interna alla macchina, che collega i container fra
loro.

\Needspace{26\baselineskip}
\subsection{Servizi interni}

Il servizio degli utenti serve solo a quello di autenticazione: non c'è motivo di renderlo raggiungibile da fuori.
Basta \textbf{non pubblicarne la porta}, togliendo \texttt{ports} (e già che ci siamo l'attributo \texttt{version}):

\begin{lstlisting}[style=yaml, name=st-microservices/docker-compose.interno.yml]
services:
 st-user-service:
  build:
    context: ./st-user-service
    dockerfile: Dockerfile
  container_name: st-user-service

 st-auth-service:
  build:
    context: ./st-auth-service
    dockerfile: Dockerfile
  container_name: st-auth-service
  ports:
    - "0.0.0.0:8081:8081"
  environment:
    - ST_USER_SERVICE_URL=http://st-user-service:8082
  depends_on:
    - st-user-service
\end{lstlisting}

\Needspace{14\baselineskip}
\begin{lstlisting}[style=shell, name=Terminale]
$ docker compose -f docker-compose.interno.yml up -d
...
 Container st-auth-service Started
$ curl -s localhost:8081/auth/janet; echo
{"username":"janet","ruolo":"studente","autorizzato":true}
$ curl -sS localhost:8082/users/janet
curl: (7) Failed to connect to localhost port 8082 after 0 ms: Couldn't connect to server
$ docker compose -f docker-compose.interno.yml ps
NAME              IMAGE                              COMMAND               SERVICE           CREATED         STATUS         PORTS
st-auth-service   st-microservices-st-auth-service   "java -jar app.jar"   st-auth-service   4 seconds ago   Up 4 seconds   0.0.0.0:8081->8081/tcp
st-user-service   st-microservices-st-user-service   "java -jar app.jar"   st-user-service   4 seconds ago   Up 4 seconds   8082/tcp
\end{lstlisting}

\Needspace{20\baselineskip}
Il servizio di autenticazione funziona come prima, perché parla con l'altro dentro la rete di Compose; dall'host, la
porta 8082 non risponde più. Nell'elenco la differenza si legge nella colonna \texttt{PORTS}: \texttt{8082/tcp} è la
porta esposta dal container (l'\texttt{EXPOSE} del Dockerfile), \texttt{0.0.0.0:8081->8081/tcp} una porta
\textbf{pubblicata} sull'host. Alla fine, \texttt{docker compose down} ferma e cancella container e rete:

\begin{lstlisting}[style=shell, name=Terminale]
$ docker compose down
time="2026-11-03T19:14:34Z" level=warning msg="[cartella del progetto]/docker-compose.yml: the attribute `version` is obsolete, it will be ignored, please remove it to avoid potential confusion"
 Container st-auth-service Stopping
 Container st-auth-service Stopped
 Container st-auth-service Removing
 Container st-auth-service Removed
 Container st-user-service Stopping
 Container st-user-service Stopped
 Container st-user-service Removing
 Container st-user-service Removed
 Network st-microservices_default Removing
 Network st-microservices_default Removed
\end{lstlisting}

(L'ultimo comando è stato eseguito con il file originale, prima della prova con il servizio interno.) I container
vengono fermati nell'ordine inverso a quello di avvio.

% ============================================================
\Needspace{40\baselineskip}
\section{Un esempio più complesso}

L'ultimo esempio delle slide è un gioco didattico sui \textit{test smell} (i difetti nel codice dei test), fatto di sei
servizi. Le immagini non vengono costruite, ma scaricate da Docker Hub, già pronte. Il file \texttt{docker-compose.yml},
ricomposto dalle due slide:

\begin{lstlisting}[style=yaml, name=docker-compose.yml (1/2)]
networks:
 testSmell-network:
  driver: bridge

services:
 userservice:
  image: mick0974/testsmellsgame:user_service
  container_name: user-service
  ports:
   - '8081:8081'
  environment:
   - SPRING_DATASOURCE_URL=jdbc:h2:file:./database/userdb
  networks:
   - testSmell-network

 compilerservice:
  image: mick0974/testsmellsgame:compiler_service
  container_name: compiler-service
  ports:
   - '8083:8083'
  networks:
   - testSmell-network

 leaderboardservice:
  image: mick0974/testsmellsgame:leaderboard_service
  container_name: leaderboard-service
  ports:
   - '8087:8087'
  environment:
   - SPRING_DATASOURCE_URL=jdbc:h2:file:./database/leaderboarddb
  networks:
   - testSmell-network
\end{lstlisting}

\Needspace{34\baselineskip}
\begin{lstlisting}[style=yaml, name=docker-compose.yml (2/2), firstnumber=34]
 exerciseservice:
  image: mick0974/testsmellsgame:exercise_service
  container_name: exercise-service
  ports:
   - '9090:9090'
  volumes:
   - assets:/usr/src/app/assets
  networks:
   - testSmell-network

 apigateway:
  image: mick0974/testsmellsgame:api_gateway
  container_name: api-gateway
  environment:
   - USER_SERVICE=http://user-service:8081
   - LEADERBOARD_SERVICE=http://leaderboard-service:8087
   - EXERCISE_SERVICE=http://exercise-service:9090
   - COMPILER_SERVICE=http://compiler-service:8083
  ports:
   - '9000:9000'
  networks:
   - testSmell-network

 frontend:
  image: mick0974/testsmellsgame:frontend
  container_name: frontend
  environment:
   - API_GATEWAY=http://api-gateway:9000
  ports:
   - '8080:8080'
  networks:
   - testSmell-network

volumes:
 assets:
  external: true
\end{lstlisting}

Le novità rispetto agli esempi precedenti:

\begin{itemize}
  \item \texttt{networks} dichiara una \textbf{rete virtuale} riservata alle connessioni fra i container,
        \texttt{testSmell-network}, e ogni servizio vi si collega;
  \item \texttt{image} indica un'immagine che non viene generata, ma è disponibile su Docker Hub
        (\texttt{utente/repository:etichetta});
  \item \texttt{environment} associa variabili d'ambiente al container: per i servizi con un database H2, il file in cui
        tenerlo; per il gateway, gli indirizzi degli altri servizi;
  \item il servizio degli esercizi espone una sua cartella come \textbf{volume}, uno spazio di memoria
        \textbf{persistente}: i file in \texttt{/usr/src/app/assets} sopravvivono alla cancellazione del container;
  \item il volume \texttt{assets} è \texttt{external}: non lo crea Compose, deve esistere già, e si può usare anche
        dall'esterno (per esempio dall'interfaccia di Docker);
  \item l'\textbf{API gateway} è l'unico servizio che il frontend contatta: instrada le richieste verso gli altri
        servizi, attraverso la rete interna (è lo schema del Capitolo 12);
  \item il \textbf{frontend} comunica con il gateway, ed è il servizio da raggiungere dall'esterno, alla porta 8080
        (della visibilità di quella porta si occuperà un altro software).
\end{itemize}

\info{Porte esposte e porte pubblicate}{
  Le slide descrivono il gateway come l'unico servizio accessibile dal frontend, e il frontend come l'unico accessibile
  dall'esterno. Nel file però \textbf{tutti} i servizi hanno una sezione \texttt{ports}, e quindi sono raggiungibili
  anche dall'host, e da lì dall'esterno. Per renderli davvero interni basta togliere \texttt{ports} dai servizi che non
  devono essere visibili, come nella sezione precedente: dentro \texttt{testSmell-network} si parlerebbero comunque.}

\Needspace{16\baselineskip}
Il \textbf{diagramma di deployment} (Capitolo 9) del sistema:

\begin{center}
\begin{tikzpicture}[
    nodo3d/.style={rectangle, draw=black!70, fill=#1!8, line width=0.7pt, font=\scriptsize, align=center,
                   minimum height=0.75cm, inner sep=3pt},
    nodo3d/.default=swblue,
    conn/.style={line width=0.7pt, draw=black!60}]
  % il nodo host
  \draw[line width=0.8pt, fill=swlight] (-0.4,-4.3) rectangle (13.4,1.3);
  \draw[line width=0.8pt, fill=swlight!60] (-0.4,1.3) -- (-0.1,1.6) -- (13.7,1.6) -- (13.4,1.3);
  \draw[line width=0.8pt, fill=swlight!40] (13.4,1.3) -- (13.7,1.6) -- (13.7,-4.0) -- (13.4,-4.3);
  \node[font=\scriptsize\bfseries, anchor=north west] at (-0.3,1.25) {Docker Host Machine};
  \node[nodo3d=swgreen] (fe) at (2.0,0.2) {frontend\\ \texttt{\guillemotleft container\guillemotright}};
  \node[nodo3d=sworange] (gw) at (6.5,0.2) {api-gateway\\ \texttt{\guillemotleft container\guillemotright}};
  \node[nodo3d] (us) at (0.9,-2.6) {user-service\\ \texttt{\guillemotleft container\guillemotright}\\ {\tiny userdb (H2)}};
  \node[nodo3d] (lb) at (4.2,-2.6) {leaderboard-service\\ \texttt{\guillemotleft container\guillemotright}\\ {\tiny leaderboarddb (H2)}};
  \node[nodo3d] (cs) at (7.6,-2.6) {compiler-service\\ \texttt{\guillemotleft container\guillemotright}};
  \node[nodo3d] (es) at (10.9,-2.6) {exercise-service\\ \texttt{\guillemotleft container\guillemotright}};
  \draw[conn] (fe) -- node[font=\tiny, above]{HTTP (API\_GATEWAY)} (gw);
  \draw[conn] (gw) -- (us); \draw[conn] (gw) -- (lb); \draw[conn] (gw) -- (cs); \draw[conn] (gw) -- (es);
  \node[font=\tiny, text=swgray, align=center] at (10.4,-0.9) {testSmell-network\\ (bridge Docker)};
  \node[font=\tiny, anchor=south] at (fe.north) {porta 8080};
  % il volume esterno
  \node[draw=swteal!80!black, fill=swteal!12, cylinder, shape border rotate=90, aspect=0.25, minimum width=1.4cm,
        minimum height=1.0cm, font=\tiny, align=center] (vol) at (12.6,-5.6) {assets\\ (volume\\ esterno)};
  \draw[conn, dashed] (es) -- (vol);
\end{tikzpicture}
\end{center}

Tutti i container girano sullo stesso nodo, la macchina che ospita Docker, collegati dalla rete \texttt{bridge};
fuori dal nodo c'è soltanto il volume esterno con le risorse degli esercizi.

% ============================================================
\Needspace{24\baselineskip}
\gbox{In sintesi}{azure-gradient-6!80!black}{
\begin{itemize}
  \item Il software moderno è fatto di molti servizi, con stack diversi, da far girare in molti ambienti: la
        \textbf{matrice infernale} delle combinazioni.
  \item Come i container del trasporto merci, i \textbf{container} del software impacchettano un servizio con tutto ciò
        che gli serve, in un formato standard che si esegue allo stesso modo ovunque.
  \item A differenza delle macchine virtuali, i container \textbf{condividono il kernel} dell'host: sono più leggeri e
        si avviano in pochi secondi, ma richiedono la stessa architettura e un kernel Linux.
  \item \textbf{Docker}: il client parla con il demone \texttt{dockerd}, che gestisce immagini e container e scarica
        le immagini da un registry. Un'\textbf{immagine} (di sola lettura, a strati) si ottiene con \texttt{pull},
        \texttt{load} o \texttt{build} da un \textbf{Dockerfile}; un container è un'immagine in esecuzione.
  \item \textbf{Docker Compose} descrive un'applicazione in un file YAML: servizi, porte pubblicate, variabili
        d'ambiente, dipendenze di avvio, politiche di riavvio, reti e volumi.
  \item I container di un progetto Compose si trovano \textbf{per nome} sulla rete del progetto; un servizio senza
        \texttt{ports} resta interno.
\end{itemize}
}
